\chapter{Αναπαραγωγή των μετρήσεων}
\label{app:reproducibility}

Το παράρτημα αυτό περιγράφει τα βήματα αναπαραγωγής των μετρήσεων της εργασίας.
Η διαδικασία περιλαμβάνει την εγκατάσταση των εργαλείων, την εκτέλεση του
κώδικα, την παραγωγή των διαγραμμάτων και της αναφοράς. Οι εντολές αφορούν το
σύστημα της ενότητας~\ref{sec:meth-machine} και εκτελούνται από τη ρίζα του
αποθετηρίου (\url{https://github.com/ioannisam/topk}).

\section{Απαιτήσεις λογισμικού}
\label{sec:repro-requirements}

Και οι τρεις υλοποιήσεις απαιτούν μεταγλωττιστή \tl{C++} που υποστηρίζει το
πρότυπο \tl{C++20}, καθώς και \tl{cmake} έκδοσης \tl{3.18} ή νεότερης. Οι
διανυσματικές εντολές δεν αποτελούν απαίτηση μεταγλώττισης. Η επιλογή ανάμεσα
σε \tl{AVX-512}, \tl{AVX2} και βαθμωτή εκτέλεση γίνεται κατά τον χρόνο
εκτέλεσης, με βάση τις δυνατότητες που αναφέρει ο επεξεργαστής. Επειδή οι
μετρήσεις έγιναν σε μηχάνημα με \tl{AVX-512}, η απουσία του οδηγεί σε
διαφορετική διαδρομή κώδικα και τα αποτελέσματα δεν είναι άμεσα συγκρίσιμα.

Η μονάδα γραφικών απαιτεί την εργαλειοθήκη \tl{CUDA}, η οποία περιλαμβάνει τον
μεταγλωττιστή \tl{nvcc} και τη βιβλιοθήκη \tl{CUB}. Οι μετρήσεις ισχύος
προϋποθέτουν επιπλέον τη βιβλιοθήκη \tl{NVML} του οδηγού.

Η μονάδα νευρωνικής επεξεργασίας απαιτεί τη στοίβα εκτέλεσης \tl{XRT} και το
πρόσθετο \tl{XDNA}. Η μεταγλώττιση των πυρήνων χρειάζεται το \tl{MLIR-AIE} σε
περιβάλλον \tl{Python}. Αυτό το περιβάλλον πρέπει να βρίσκεται στη ρίζα του
αποθετηρίου με το όνομα \tl{npu\_venv}, καθώς εκεί αναζητείται από τα σενάρια
εκτέλεσης. Ο χρήστης πρέπει να ανήκει στις ομάδες πρόσβασης της συσκευής και να
έχει επαρκές όριο κλειδωμένης μνήμης, διαφορετικά η συσκευή δεν αναγνωρίζεται.
Η μεταβλητή \tl{XILINX\_XRT} πρέπει να δείχνει στην εγκατάσταση της \tl{XRT}.

Η μέτρηση ενέργειας διαβάζει τους μετρητές \tl{RAPL} από το
\tl{/sys/class/powercap}, διαδικασία που απαιτεί δικαιώματα διαχειριστή. Σε
μεμονωμένη εκτέλεση χωρίς αυτά, η μέτρηση δηλώνεται ρητά μη διαθέσιμη αντί να
καταγράψει μηδενικά (ενότητα~\ref{sec:meth-energy}). Αντίθετα, η
αυτοματοποιημένη σάρωση ακυρώνεται εντελώς για να αποφευχθεί η παραγωγή κενού
συνόλου δεδομένων. Τέλος, η υποδομή ανάλυσης απαιτεί τις βιβλιοθήκες
\tl{Python} του αποθετηρίου. Το αντίστοιχο περιβάλλον δημιουργείται αυτόματα
κατά την πρώτη εκτέλεση.

\section{Μεταγλώττιση}
\label{sec:repro-build}

Η μεταγλώττιση όλων των υλοποιήσεων γίνεται με μία εντολή:

\selectlanguage{english}
\begin{verbatim}
make build-all
\end{verbatim}
\selectlanguage{greek}

Κάθε αρχιτεκτονική μπορεί να μεταγλωττιστεί και αυτόνομα, με τους στόχους
\tl{build-cpu}, \tl{build-gpu} και \tl{build-npu}, κάτι που είναι χρήσιμο σε
σύστημα που δεν διαθέτει και τις τρεις μονάδες. Τα παραγόμενα εκτελέσιμα
τοποθετούνται στον κατάλογο \tl{build}, δύο ανά αρχιτεκτονική: ένα για τις
μετρήσεις \tl{top-k} και ένα για τις μετρήσεις των ταβανιών της
ενότητας~\ref{sec:meth-roofline}.

Η μεταγλώττιση για τη μονάδα νευρωνικής επεξεργασίας παράγει επιπλέον τα αρχεία
διαμόρφωσης του υλικού. Επειδή η απουσία τους οδηγεί σε σιωπηλή παράλειψη των
αντίστοιχων μετρήσεων, η παρουσία τους ελέγχεται ρητά:

\selectlanguage{english}
\begin{verbatim}
make verify-artifacts
\end{verbatim}
\selectlanguage{greek}

Ο έλεγχος εξετάζει τα δύο αρχεία διαμόρφωσης και τα τρία εκτελέσιμα \tl{top-k}.
Προϋποθέτει πλήρη μεταγλώττιση, οπότε αποτυγχάνει αυτόματα αν έχει
μεταγλωττιστεί μόνο μία αρχιτεκτονική.

\section{Εκτέλεση των πειραμάτων}
\label{sec:repro-run}

Ο συνιστώμενος τρόπος αναπαραγωγής είναι η συνολική εκτέλεση της διαδικασίας με
μία εντολή:

\selectlanguage{english}
\begin{verbatim}
make benchmark GPU_W=50
\end{verbatim}
\selectlanguage{greek}

Αυτή η εντολή διασφαλίζει τη σωστή αλληλουχία βημάτων. Αν οι συνθήκες
καθηλωθούν εκ των υστέρων ή χρησιμοποιηθούν παλαιά εκτελέσιμα, τα αποτελέσματα
χάνουν τη συγκρισιμότητά τους. Η εντολή επιβάλλει την κατάλληλη σειρά και
αποτρέπει τις παραλείψεις.

Η ακολουθία περιλαμβάνει έξι στάδια. Αρχικά διαγράφονται τα παλαιά αποτελέσματα
και εκτελέσιμα. Ακολουθεί η μεταγλώττιση και ο έλεγχος των αρχείων
(ενότητα~\ref{sec:repro-build}). Στο τρίτο στάδιο καθηλώνονται οι συνθήκες του
συστήματος και παρέχεται πρόσβαση στους μετρητές ενέργειας. Ακολουθούν οι
τέσσερις σαρώσεις μετρήσεων, η παραγωγή των διαγραμμάτων και, τελικά, η
αποκατάσταση του συστήματος.

Η διαδικασία ελέγχει πρώτα τη διαθεσιμότητα του \tl{sudo} και του περιβάλλοντος
\tl{npu\_venv}. Αν κάτι από τα δύο απουσιάζει, η εντολή τερματίζει άμεσα, χωρίς
να αλλοιώσει το σύστημα ή τα δεδομένα.

Τα δικαιώματα διαχειριστή ζητούνται μία φορά και παραμένουν ενεργά. Οι συνθήκες
αποκαθίστανται πάντα στο τέλος, ακόμη και αν υπάρξει σφάλμα ή διακοπή από τον
χρήστη, αποτρέποντας τον εγκλωβισμό του συστήματος σε καθηλωμένη κατάσταση. Η
αποτυχία μιας μεμονωμένης σάρωσης δεν σταματά τις υπόλοιπες, καταγράφεται απλώς
ως προειδοποίηση.

Το όρισμα \tl{GPU\_W} ορίζει το όριο ισχύος της μονάδας γραφικών, το οποίο
πρέπει να παραμείνει σταθερό (ενότητα~\ref{sec:meth-conditions}). Η διαδικασία
διαγράφει αμέσως τα υπάρχοντα αποτελέσματα, συνεπώς ο χρήστης οφείλει πρώτα να
κρατήσει αντίγραφα όσων δεδομένων θέλει να διατηρήσει.

\subsection{Εκτέλεση επιμέρους βημάτων}
\label{sec:repro-steps}

Για μερική ή επαναληπτική εκτέλεση, κάθε στάδιο είναι διαθέσιμο και ξεχωριστά.
Οι συνθήκες του συστήματος ελέγχονται με:

\selectlanguage{english}
\begin{verbatim}
make pin ARGS=50      # governor, boost off, GPU cap 50 W
make pin-show         # report the current state
make unpin            # restore governor, boost, GPU cap
\end{verbatim}
\selectlanguage{greek}

Οι τέσσερις ανεξάρτητες σαρώσεις καταγράφουν τα ακατέργαστα αποτελέσματά τους
σε χωριστά αρχεία:

\selectlanguage{english}
\begin{verbatim}
make measure-cases      # timing over the parameter space
make measure-energy     # energy, requires RAPL access
make measure-dists      # input-distribution sensitivity
make measure-roofline   # bandwidth, compare-exchange roofs
\end{verbatim}
\selectlanguage{greek}

Οι στόχοι δέχονται πρόσθετα ορίσματα μέσω της μεταβλητής \tl{ARGS},
περιορίζοντας τον παραμετρικό χώρο. Αυτό διευκολύνει τον γρήγορο έλεγχο της
εγκατάστασης. Στη χειροκίνητη εκτέλεση, ο χρήστης αναλαμβάνει την ευθύνη για τη
σωστή σειρά και τη διατήρηση σταθερών συνθηκών.

\subsection{Παράμετροι εκτέλεσης}
\label{sec:repro-cli}

Ο πίνακας~\ref{tab:cli} συνοψίζει τις παραμέτρους κάθε εκτελέσιμου, τις
επιτρεπτές τιμές, τις προεπιλογές και τις τιμές των σαρώσεων.

Συχνά οι προεπιλογές διαφέρουν από τις τελικές τιμές των σαρώσεων. Το εύρος
τιμών ορίστηκε έως $10^9$ αντί της προεπιλογής, για να αποφευχθεί η υπερβολική
συσσώρευση ισοβαθμιών σε εισόδους εκατομμυρίων στοιχείων. Οι πολλές ισοβαθμίες
θα ευνοούσαν τεχνητά τους αλγορίθμους που απορρίπτουν γρήγορα στοιχεία.

Το εύρος αυτό αφορά τους τέσσερις από τους πέντε τύπους. Για τον \tl{half}, το
εκτελέσιμο περικόπτει τα άκρα στο $\pm 65504$, τη μέγιστη πεπερασμένη τιμή του
τύπου. Σε αυτή την περίπτωση οι ισοβαθμίες παραμένουν αναπόφευκτες, λόγω του
στενού εύρους και της αδυναμίας του \tl{half} να αναπαραστήσει όλους τους
ακεραίους πάνω από το $2048$. Η ιδιότητα αυτή πηγάζει από τον ίδιο τον τύπο και
επηρεάζει άμεσα την ερμηνεία των αποτελεσμάτων.

Ο σπόρος τυχαίων αριθμών παράγεται ντετερμινιστικά από τα \tl{q} και \tl{k} της
κάθε περίπτωσης. Έτσι διασφαλίζεται η επαναληψιμότητα χωρίς να δοθεί ποτέ η
ίδια είσοδος σε διαφορετικές συνθήκες.

Η κύρια σάρωση χρησιμοποιεί τις τιμές της τέταρτης στήλης. Οι ειδικές σαρώσεις
περιορίζουν τις παραμέτρους: η σάρωση κατανομών εξετάζει \tl{q} από 20 έως 24,
τύπους \tl{float} και \tl{half}, και άκρα \tl{k} 8 και 131072, με δύο σπόρους
ανά περίπτωση. Η σάρωση ενέργειας εξετάζει \tl{q} από 16 έως 24 με τρεις
επαναλήψεις, παραλείποντας τον έλεγχο ορθότητας για να μην επιβαρυνθεί άσκοπα η
καταναλισκόμενη ισχύς.

\begin{table}[H]
    \centering
    \small
    \begin{tabular}{@{}lp{2.9cm}lp{2.6cm}p{4.0cm}@{}}
        \hline
        Παράμετρος & Τιμές & Προεπιλογή & Στις σαρώσεις & Σημασία \\
        \hline
        \tl{q} & ακέραιος \tl{0}--\tl{63} & --- & \tl{3}--\tl{24} & Μέγεθος εισόδου $n = 2^q$· είναι η μόνη υποχρεωτική \\
        \tl{k} & θετικός ακέραιος & \tl{n} & \tl{8}, \tl{256}, \tl{4096}, \tl{131072} & Πλήθος ζητούμενων στοιχείων· περιορίζεται στο \tl{n} \\
        \tl{mode} & \tl{min} $|$ \tl{max} & \tl{max} & \tl{max} & Επιλογή μικρότερων ή μεγαλύτερων \\
        \tl{dtype} & \tl{int} $|$ \tl{uint} $|$ \tl{float} $|$ \tl{double} $|$ \tl{half} & \tl{int} & και οι πέντε & Τύπος στοιχείου \\
        \tl{algo} & \tl{bitonic} $|$ \tl{map\_reduce} $|$ \tl{gt} & \tl{bitonic} & και οι τρεις & Αλγόριθμος προς εκτέλεση \\
        \tl{run} & \tl{full} $|$ \tl{trunc} $|$ \tl{both} & \tl{trunc} & \tl{trunc} & Παραλλαγή του διτονικού δικτύου \\
        \tl{seed} & μη αρνητικός ακέραιος & \tl{42} & παράγεται από τα \tl{q}, \tl{k} & Σπόρος της γεννήτριας τυχαίων αριθμών \\
        \tl{dist} & \mbox{\tl{uniform}} $|$ \mbox{\tl{normal}} $|$ \mbox{\tl{trimodal}} $|$ \mbox{\tl{sorted}} $|$ \mbox{\tl{reverse}} $|$ \mbox{\tl{adversarial}} & \tl{uniform} & \tl{uniform}· πέντε στη σάρωση κατανομών & Κατανομή της εισόδου \\
        \tl{min} / \tl{max} & ακέραιοι & \tl{0} / \tl{1000} & \tl{0} / $10^9$ & Εύρος τιμών· απαιτείται \tl{max} $\geq$ \tl{min}· περικόπτεται στο εύρος του τύπου \\
        \tl{threads} & θετικός ακέραιος & \tl{0} (αυτόματο) & \tl{8} & Πλήθος νημάτων· η \tl{NPU} το αγνοεί \\
        \tl{verify} & \tl{true} $|$ \tl{false} & \tl{false} & \tl{true} & Έλεγχος ως προς την υλοποίηση επαλήθευσης \\
        \tl{debug} & \tl{true} $|$ \tl{false} & \tl{true} & \tl{false} & Αναλυτική έξοδος διαγνωστικών \\
        \hline
    \end{tabular}
    \caption{Οι παράμετροι εκτέλεσης και οι προεπιλογές τους.}
    \label{tab:cli}
\end{table}

\section{Παραγωγή των διαγραμμάτων}
\label{sec:repro-plots}

Η υποδομή ανάλυσης διαβάζει τα ακατέργαστα αποτελέσματα, τα μετατρέπει σε
πίνακες \tl{CSV} και παράγει τα διαγράμματα:

\selectlanguage{english}
\begin{verbatim}
make plot ARGS="--plot all"
\end{verbatim}
\selectlanguage{greek}

Το όρισμα \tl{--plot} επιλέγει τα διαγράμματα και το \tl{--input} τη σάρωση (η
προεπιλογή είναι η σάρωση χρόνου). Επειδή τα υπόλοιπα διαγράμματα χρησιμοποιούν
δικά τους αρχεία, συνιστώνται οι συνδυαστικοί στόχοι, οι οποίοι εκτελούν άμεσα
τη μέτρηση και την παραγωγή:

\selectlanguage{english}
\begin{verbatim}
make benchmark-cases
make benchmark-energy
make benchmark-dists
make benchmark-roofline
\end{verbatim}
\selectlanguage{greek}

Τα αρχεία οργανώνονται ανά τύπο δεδομένων και κατηγορία. Διαγράμματα που
αφορούν αποκλειστικά το υλικό (π.χ. κλίμακα εύρους ζώνης) αποθηκεύονται σε
ξεχωριστό κατάλογο, καθώς δεν εξαρτώνται από τον τύπο.

Δύο παράμετροι καθορίζουν την απεικόνιση. Η \tl{--energy-metric} επιλέγει
μεταξύ μεικτής ή καθαρής ενέργειας. Οι σαρώσεις της εργασίας χρησιμοποιούν την
καθαρή (ενότητα~\ref{sec:meth-energy}). Η παράμετρος \tl{--error-bars} ορίζει
τη ζώνη διασποράς (τυπική απόκλιση, δεκατημόρια ή καμία). Η ζώνη σχεδιάζεται
μόνο αν υπάρχουν τουλάχιστον δύο δείγματα ανά σημείο, δηλαδή στη σάρωση
κατανομών (δύο σπόροι) και ενέργειας (τρεις επαναλήψεις). Στην αυτοματοποιημένη
εκτέλεση, οι ζώνες απενεργοποιούνται πλήρως.
